\section{Cyclotomic divisibility tools}\label{cyclo-tools}

In this section, we develop the tools needed to prove our results on cyclotomic factor decompositions and vanishing sums of roots of unity. Some of the notation here has been borrowed from \cite{LaLo1} and adapted to our setting.

\subsection{Multisets}
 
We will work in the ambient group $\ZZ_M$, where $M=p_1^{n_1}\dots p_K^{n_K}$, $p_1,\dots,p_K$ are distinct primes, and $n_1,\dots,n_K\in\NN$. We will use $m$ and $N$ (possibly with subscripts) to denote divisors of $M$. For the purpose of proving Theorem \ref{main theorem}, it would be sufficient to consider the case $K=2$. The discussion in Section \ref{sec-examples} will require the more general notation. 

We use $A(X)$, $B(X)$, etc. to denote polynomials modulo $X^M-1$ with integer coefficients. 
Each such polynomial $A(X)=\sum_{a\in\ZZ_M} w_A(a) X^a$ is associated with a weighted multiset in $\ZZ_M$, which we will also denote by $A$, with weights $w_A(x)$ assigned to each $x\in\ZZ_M$. (If the coefficient of $X^x$ in $A(X)$ is 0, we set $w_A(x)=0$.) In particular, if $A$ has $\{0,1\}$ coefficients, then
$w_A$ is the characteristic function of a set $A\subset \ZZ_M$. We will use $\calm(\ZZ_M)$ to denote the 
family of all weighted multisets in $\ZZ_M$, and reserve the notation $A\subset \ZZ_M$ for sets.
We will also use $\calm^+$ to denote the family of all weighted multisets in $\ZZ_M$ with nonnegative weights:
$$
\calm^+=\{A\in\calm(\ZZ_M): \ w_A(a)\geq 0\hbox{ for all }a\in\ZZ_M\}.
$$
Let $A\in\calm^+(\ZZ_M)$, with the corresponding mask polynomial $A(X)$. We use $|A|$ to denote the cardinality of $A$ with multiplicity, so that $|A|=A(1)=\sum_{x\in\ZZ_M}w_A(x)$.
If $\Lambda\subset\ZZ_M$, we use $A\cap\Lambda$ to denote intersection with multiplicity, so that $w_{A\cap\Lambda}(x)=w_A(x)w_\Lambda(x)$.


If $N|M$, then any $A\in \calm(\ZZ_M)$ induces a weighted multiset $A$ mod $N$ in $\ZZ_N$, with the corresponding mask polynomial $A(X)$ mod $(X^N-1)$ and induced weights 
\begin{equation}\label{induced-weights}
w_A^N(x)= \sum_{x'\in\ZZ_M: x'\equiv x\,{\rm mod}\, N} w_A(x'),\ \ x\in\ZZ_N.
\end{equation}
We extend the multiset notation to $\ZZ_N$, so that for example $\calm(\ZZ_N)$ and $\calm^+(\ZZ_N)$ denote the appropriate families of multisets.

We use convolution notation $A*B$ to denote the weighted sumset of $A,B\in\calm(\ZZ_M)$:
 $$
 (A*B)(X)=A(X)B(X),\ \ 
 w_{A*B}(x)=(w_A*w_B)(x)=\sum_{y\in\ZZ_M} w_A(x-y)w_B(y).$$
 If one of the sets is a singleton, say $A=\{x\}$, we write $x*B=\{x\}*B$.  

\subsection{Grids and fibers}\label{grids-fibers}
We encourage the reader to use the geometrical interpretation from \cite{LaLo1}, based on the Chinese Remainder Theorem. Specifically, let $M_i=M/p_i^{n_i}$ for $i=1,\dots,K$. Then any $x\in\ZZ_M$ can be written uniquely as 
$$
x\equiv \sum_{i=1}^{K} x_iM_i\hbox{ mod }M,\hbox{ where }x_i\in\{0,1,\dots,p_i-1\}.
$$
Thus $\ZZ_M$ can be thought of as a $K$-dimensional lattice $\ZZ_{p_1^{n_1}}\oplus\dots\oplus \ZZ_{p_K^{n_K}}$. In this interpretation, $(x_1,\dots,x_K)$ are the coordinates of $x\in\ZZ_M$, and each of the primes $p_1,\dots,p_K$ corresponds to one of the cardinal directions. A similar picture, possibly with fewer directions, applies to $\ZZ_N$ with $N|M$. 

For $D|N|M$, a {\em $D$-grid} in $\ZZ_N$ is a set of the form
$$
\Lambda^N(x,D):= x*D\ZZ_N=\{x'\in\ZZ_N:\ D|(x-x')\}
$$
for some $x\in\ZZ_N$.  
If $N=\prod_{j=1}^K p_j^{\alpha_j}$ is a divisor of $M$, with  $0\leq \alpha_j\leq n_j$, we let 
$$
D(N):= \prod_{j=1}^K p_j^{\gamma_j},\ \hbox{ where }
\gamma_j=\max(0,\alpha_j-1)\hbox{ for }j=1,\dots,K.
$$  


Let $p_i$ be a prime factor of $N$. An {\em $N$-fiber in the $p_i$ direction} is a set of the form $x*F^N_i\subset\ZZ_N$, where $x\in\ZZ_N$ and
\begin{equation}\label{def-Fi}
F^N_i=\{0,N/p_i,2N/p_i,\dots,(p_i-1)N/p_i\}.
\end{equation}
Thus $x*F^N_i=\Lambda^N(x,N/p_i)$. 





%

\subsection{Cuboids and structure results}
As before, we work in $\ZZ_M$, where $M=\prod_{i=1}^K p_i^{n_i}$, and let $A\in\calm(\ZZ_M)$.
We will use the following notation from \cite{LaLo1}. For multisets $\Delta\in\calm(\ZZ_N)$, where $N|M$, we define the \textbf{$\Delta$-evaluations of $A$ in $\ZZ_N$:}
\begin{equation}\label{delta-eval}
\bbA^N[\Delta]=\sum_{x\in\ZZ_N}w_A^N(x)w_\Delta^N(x).
\end{equation}
The following special case is of particular interest.

\begin{definition}
\label{def-N-cuboids}
Let $M$ and $N$ be as above, so that $N=\prod_{i=1}^K p_i^{n_i-\alpha_i}$, with $0\leq \alpha_i\leq n_i$ for each $i=1,\dots,K$. An \textbf{$N$-cuboid} is a multiset $\Delta \in \calm (\zz_N)$ associated to a mask polynomial of the form
\begin{equation}\label{def-delta}
\Delta(X)= X^c\prod_{j:p_j|N} (1-X^{d_jN/p_j})
\end{equation}
 with $(d_j,p_j)=1$ for all $j$. 
\end{definition} 
 


The geometric interpretation of $N$-cuboids $\Delta$ is as follows. With notation as in Definition \ref{def-N-cuboids}, 
recall that $D(N)=N/\prod_{j:p_j|N}p_j$. Then the ``vertices'' $x\in\ZZ_N$ with $ w^N_\Delta(x)\neq 0$ form a full-dimensional rectangular box in the grid $\Lambda(c,D(N))$, with one vertex at $c$ and alternating $\pm 1$ weights. 



The following cyclotomic divisibility test has been known and used previously in the literature. The equivalence between (i) and (iii) is the Bruijn-R\'edei-Schoenberg theorem on the structure of vanishing sums of roots of unity (see \cite{deB}, \cite{LL}, \cite{Mann}, \cite{Re1}, \cite{Re2}, \cite{schoen}). For the equivalence (i) $\Leftrightarrow$ (ii), see e.g.  \cite[Section 3]{Steinberger}, \cite[Section 3]{KMSV}.


\begin{proposition}\label{cuboid}
Let $A\in\calm(\ZZ_N)$. Then the following are equivalent:

\medskip

(i) $\Phi_N(X)|A(X)$,


\medskip

(ii) For all $N$-cuboids $\Delta$, we have
\begin{equation}\label{id-3a}
\bbA^N[\Delta]=0,
\end{equation}

\medskip

(iii) $A$ mod $N$ is a linear combination of $N$-fibers, so that
$$A(X)=\sum_{i:p_i|N} P_i(X)F^N_i(X) \mod X^N-1,$$
where $P_i(X)$ have integer (but not necessarily nonnegative) coefficients.
\end{proposition}

Proposition \ref{cuboid} can be strengthened as follows if $N$ has only two distinct prime factors. This goes back to the work of de Bruijn \cite{deB}; a self-contained proof is provided in \cite[Theorem 3.3]{LL}.

\begin{lemma}\label{structure-thm} 
Let $A\in\calm^+(\ZZ_N)$.
Assume that $\Phi_{N}|A$, where $N$ has two distinct prime factors $p_1,p_2$. Then $A$ mod $N$ is a linear combination of $N$-fibers with nonnegative weights. In other words,
$$A(X)=P_1(X)F^N_1(X)+ P_2(X)F^N_2(X) \mod X^N-1,$$
where $P_1,P_2$ are polynomials with nonnegative coefficients.

\end{lemma}



\begin{lemma}\label{grid-split}
Assume that $N|M$ and $A\in\calm(\ZZ_N)$. Let $m|D(N)$.Then $\Phi_N|A$ if and only if $\Phi_N|(A\cap\Lambda)$ for every $m$-grid $\Lambda$.
\end{lemma}

\begin{proof}
If $\Delta$ is an $N$-cuboid with one vertex $z\in\ZZ_N$, then all its vertices are contained in $\Lambda(x,D(N))$. Since $m|D(N)$, it follows that any $m$-grid $\Lambda$ containing any vertex of $\Delta$ must contain all of its vertices. The lemma now 
follows from the equivalence (i) $\Leftrightarrow$ (ii) in Proposition \ref{cuboid}.
\end{proof}

%


\section{Proof of Theorem \ref{main theorem}}\label{main-proof}


We are now ready to prove Theorem \ref{main theorem}. For the reader's convenience, we state it here again in the notation of Section \ref{cyclo-tools}.


\begin{proposition}\label{2primebound}
Assume that $K=2$, and write $p=p_1,q=p_2$ for short.
Let $A\in\calm^+(\ZZ_M)$, and let $m_1,m_2,\dots,m_r$ be divisors of $M$ such that $m_j=p^{\alpha_j} q^{\beta_j}$, where
$$
1\leq\alpha_1<\alpha_2<\dots <\alpha_r.
$$
Assume that $\Phi_{m_1}\Phi_{m_2}\dots \Phi_{m_r}|A$, and that $q\nmid |A|$.
Then $|A|\geq p^r$.
\end{proposition}

\begin{proof}
We proceed by induction in $r$. For the base case, suppose that $r=1$. By Lemma \ref{structure-thm}, $A$ mod $m_1$ is a union of $m_1$-fibers in the $p$ and $q$ directions. Since $q\nmid |A|$, at least one of these fibers must be in the $p$ direction. Hence $|A|\geq p$. 

Suppose now that $r\geq 2$, and that the proposition is true with $r$ replaced by $r-1$. 
Let $D:=D(p^{\alpha_1})=p^{\alpha_1-1}$. 
We write $\ZZ_M$ as a disjoint union of grids $\Lambda(y_i,D)$, where $y_0:=0,y_1,\dots,y_{D-1}\in\ZZ_M$. 
Let $A_i=A\cap \Lambda(y_i,D)$.
By Lemma \ref{grid-split}, we have $\Phi_{m_1}\Phi_{m_2}\dots \Phi_{m_r}|A_i$ for each $i$. Moreover, since $q\nmid |A|$, there exists at least one $i$ such that $q\nmid |A_i|$. Without loss of generality, we may assume that $q\nmid|A_0|$, We will prove that $|A_0|\geq p^r$.


Write $m=m_1$ for short.
Applying Lemma \ref{structure-thm} to $A_0$ on the scale $m$, we see that $A_0$ mod $m$ is a linear combination of $m$-fibers in the $p$ and $q$ directions with nonnegative coefficients. Taking into account that $A_0\subset D\ZZ_M$, we see that
\begin{equation}\label{A0-decomp}
A_0(X)\equiv P_1(X^D)F^{m}_1(X)+ P_2(X^D)F^{m}_2(X) \mod X^{m}-1,
\end{equation}
where $P_1,P_2$ are polynomials with nonnegative coefficients, $F^{m}_1(X)=1+X^{m/p}+\dots +X^{(p-1)m/p}$ is an $m$-fiber in the $p=p_1$ direction,
and $F^{m}_2=1+X^{m/q}+\dots+X^{(q-1)m/q}$ is an $m$-fiber in the $q=p_2$ direction. 

Moreover, we have the following simplification. For each $a\in\ZZ_M$, we may use the Chinese Remainder Theorem to write 
\begin{equation}\label{decomp-a}
aD\equiv a_1p^{\alpha_1} + a_2m/p\mod m,
\end{equation}
 where $a_1,a_2\in\ZZ_M$. Then
\begin{align*}
X^{aD}F^m_1(X)&=X^{a_1p^{\alpha_1}}X^{a_2m/p}
\Big(1+X^{m/p}+\dots +X^{(p-1)m/p}\Big)
\\
&= X^{a_1p^{\alpha_1}}\Big(X^{a_2m/p}+X^{(a_2+1)m/p}+\dots +X^{(a_2+p-1)m/p}\Big)
\\
&\equiv X^{a_1p^{\alpha_1}}F^m_1(X) \mod X^m-1.
\end{align*}
Applying this to every monomial in $P_1(X^D)$ if necessary, we may assume that $P_1$ satisfies
$$
P_1(X^D)=P_0(X^{p^{\alpha_1}}),
$$
where $P_0(X)$ is a polynomial with nonnegative coefficients.
Since $q\nmid |A_0|$ and $|F^{m}_2|=q$, it follows that 
\begin{equation}\label{qndiv P1}
q\nmid P_1(1).
\end{equation}

We now split up $A_0$ further, as follows. 
Let $x_0:=0,x_1,\dots,x_{p-1}$ be points in $\Lambda(0, D)$ such that $(x_i-x_j,M)=M/p^{\alpha_1-1}$ for $i\neq j$.
Using the Chinese Remainder Theorem as in (\ref{decomp-a}), we write 
$\Lambda(0, D)=\bigcup_{j=0}^{p-1}\Lambda_j$, where $\Lambda_j=\Lambda(jm/p,pD)$. Let
$A_{0,j}=A_0\cap\Lambda_j$.

The geometric idea in the next step is as follows. Think of $p$ and $q$ as two directions in a plane. We may then interpret $\{\Lambda_j\}$ as a decomposition of a 2-dimensional grid $\Lambda(0,D)$ into a system of parallel lines, each perpendicular to the $p$ direction. 
Consider the decomposition (\ref{A0-decomp}). The fibers in the $p$ direction are ``orthogonal'' to the parallel lines, so each such fiber contributes one point to each line. The fibers in the $q$ direction are parallel to $\Lambda_j$, hence any translated copy of $F^{m}_2$ is either contained fully in $\Lambda_j$ or disjoint from it.

We now write this out more explicitly. For the first part of (\ref{A0-decomp}), we have
$$
P_1(X^D)F^{m}_1(X)= P_0(X^{p^{\alpha_1}}) \Big(1+X^{m/p}+\dots +X^{(p-1)m/p}\Big),
$$
and for any monomial $cX^{ap^{\alpha_1}}$ appearing in $P_0$, we have 
$ap^{\alpha_1}+ jm/p \in \Lambda_j$ for $j=0,1,\dots,p-1$. For the second part, we consider the multiset associated to $P_2(X^D)$, decompose it into disjoint multisets in $\calm(\Lambda_j)$ for $j\in\{0,1,\dots,p-1\}$, and use that $F^m_2\subset \Lambda_0$. Hence for each $0 \leq j \leq p - 1$, there are polynomials $Q_j$ such that
$$
A_{0,j}(X)=P_1(X) X^{x_j} + Q_j (X) F^{m}_2(X) \mod X^{m}-1.
$$
This decomposition is analogous to that of $A_0$ in \eqref{A0-decomp}.
Observe that
$$
|A_{0,j}|=A_{0,j}(1)=P_1(1) + Q_j(1)F^m_2(1)= P_1(1) + Q_j(1)q.
$$
By (\ref{qndiv P1}), $|A_{0,j}|$ is not divisible by $q$.
On the other hand, by another application of Lemma \ref{grid-split}, we have $\Phi_{m_2}\dots \Phi_{m_r}|A_{0,j}$ for each $j$. 

Applying the inductive assumption to $A_{0,j}$ for each $j$, we see that $|A_{0,j}|\geq p^{r-1}$. Hence
$$
|A_0|=\sum_{j=0}^{p-1}|A_{0,j}|\geq p^r,
$$
as claimed.
\end{proof}
